\section{The exact existence statement}

The set is the full finite-dimensional bipartite tensor-product quantum
set in the coordinates
\[
 q=(a_0,a_1,b_0,b_1,E_{00},E_{01},E_{10},E_{11})\in\R^8.
\]
Writing outcomes as \(a,b\in\{-1,1\}\), its probabilities are
\[
 p(a,b\mid x,y)=\frac{1+a\,a_x+b\,b_y+ab\,E_{xy}}4.
\]
Let \(P\) be the cone of affine functions nonnegative on \(\Q\).
For a real closed extension \(K\supseteq\R\), a subscript \(K\) means
base extension of the fixed semialgebraic description.

For arbitrary separately labelled real closed extensions \(K_c,K_t\),
put
\[
 \T=K_c\otimes_{\R}K_t,\qquad
 [f,q]=f_0\otimes1+\sum_{\nu=1}^{8}f_\nu\otimes q_\nu.
\]
The source coefficients belong to \(K_c\); the independent testing point
belongs to \(K_t\). These are not Alice/Bob labels. The factors may be
two labelled copies of the same field. No ordering, field property, or
algebraic-disjointness property is assumed for \(\T\).

The October manuscript's universal claim is
\begin{equation}
 \boxed{\ 
 [f,q]\in\Sigma(K_c\otimes_{\R}K_t)^2
 \quad\text{for every }K_c,K_t,\ f\in P_{K_c},\ q\in(\Q)_{K_t}.
 \ }
 \label{eq:universal}
\end{equation}
Every certificate is finite. A common number of roots and a rank-two
bound are not part of this qualitative statement.

Scheiderer's Corollary 3.19 characterizes closed convex semialgebraic
spectrahedral shadows by this tensor-SOS condition for two copies of
every real closed extension. Its ambient dimension is arbitrary; the
planar main theorem in the paper's title has a different scope.
Thus \eqref{eq:universal}, with the finite physical model below, gives
finite \(N,m\) and real symmetric matrices satisfying
\begin{equation}
 \boxed{\ 
 \Q=\left\{q:\exists y\in\R^m,\quad
 A_0+\sum_{\nu=1}^{8}q_\nu A_\nu+
       \sum_{j=1}^{m}y_jB_j\succeq0\right\}.
 \ }
 \label{eq:lift}
\end{equation}
The same paper's Theorem 3.10 adds a uniform bound on tensor-root rank
and controls semidefinite block size. That is the stronger criterion
used by the August route. \src{S1}

The external implication is correct. The Q222-specific burden is the
universal hypothesis, not evaluating the eventual matrices.
Conversely, ordinary scalar nonnegativity alone does not supply that
hypothesis.

\section{The finite physical model}

Put \(\mathcal Z=S^1\times S^1\times S^3\).
For \(z=((c_A,s_A),(c_B,s_B),\psi)\), let \(\Phi(z)\) be the eight
expectations in a real unit state \(\psi\in\R^4\), using
\[
 A_0=B_0=Z,\quad A_1=c_AZ+s_AX,\quad B_1=c_BZ+s_BX,
\]
\[
 Z=\begin{pmatrix}1&0\\0&-1\end{pmatrix},\qquad
 X=\begin{pmatrix}0&1\\1&0\end{pmatrix}.
\]

\begin{proposition}[Exact real-qubit reduction]
\[
 \Q=\operatorname{conv}\Phi(\mathcal Z)
 =\left\{\sum_{i=1}^{9}w_i\Phi(z_i):
 z_i\in\mathcal Z,\ w_i\ge0,\ \sum_iw_i=1\right\}.
\]
Consequently \(\Q\) is compact, convex and semialgebraic.
\end{proposition}

\begin{proof}
A binary observable is a Hermitian contraction \(A\), with involutive
dilation
\[
 \widehat A=
 \begin{pmatrix}A&\sqrt{I-A^2}\\
 \sqrt{I-A^2}&-A\end{pmatrix}.
\]
The same embedding \(h\mapsto(h,0)\) works for both local settings.
Applying both local embeddings preserves joint expectations.

For two involutions \(U,V\), \(C=UV\) is unitary and \(UCU=C^*\).
Pairing conjugate eigenspaces gives invariant blocks of dimension two;
real eigenvalues give one-dimensional blocks. Local basis changes put
each nontrivial pair into the real form above. One-dimensional
strategies are represented on common eigenvectors with commuting
measurement choices.

Block-diagonal measurements retain only the diagonal state blocks,
whose traces are convex weights. A complex pure vector \(x+iy\)
against a real symmetric observable \(M\) satisfies
\[
 (x+iy)^*M(x+iy)=x^TMx+y^TMy.
\]
Separate normalization gives real pure-state atoms. Mixed states are
included by pure-state decomposition. Conversely, finite mixtures are
realized by local direct sums and a state on matching block pairs.

Carath\'eodory elimination in \(\R^8\) leaves at most nine atoms.
The displayed parameter space with simplex weights is compact and
semialgebraic, and the expectation map is polynomial. Its image is
compact and semialgebraic; direct sums establish convexity.
\end{proof}

This covers complex and mixed states and arbitrary finite original
dimensions. The qubit reduction agrees with Masanes's extremal
realization theorem. \src{RFA; S3}
The fixed nine-atom formula defines \((\Q)_K\), and
\[
 f\in P_K\quad\Longleftrightarrow\quad
 \forall q\in(\Q)_K,\ f(q)\ge0
\]
is a first-order ordered-field formula. Transfer does not assert
order-topological compactness of a non-Archimedean field.

\section{The algebra that makes the proof computer-free}

\subsection{Why the two tensor factors must remain separate}

In a real closure \(K\) of \(\R(t)\), ordered with \(t>0\)
infinitesimal, the tensor
\(\Delta=t\otimes1-1\otimes t\) is nonzero, although multiplication
sends it to zero. Extend \(d/dt\) to the algebraic extension and apply
\(a\otimes b\mapsto a\,db/dt\); the image is \(-1\).

A mixed tensor with positive residue need not be a unit.
For example \(1+\Delta^2\) is not invertible in \(K\otimes_\R K\).
Embed the second copy of \(\R(t)\) into \(K(i)\) by \(t\mapsto t+i\),
and extend to an embedding of its algebraic real closure. The induced
homomorphism sends \(\Delta\) to \(-i\), so it sends \(1+\Delta^2\)
to zero. A unit cannot have zero image in a field.

Thus a source-field Schur complement can prove source geometry without
automatically authorizing the same inverse in the tensor ring.
The proof needs pure-factor operations or exact finite tensor identities.
\src{TF}

\subsection{Legal restoration, matrix pairing and mixture identities}

For \(s>0\) in \(K_c\) and \(t>0\) in \(K_t\),
\[
 (s\otimes t)F\in\sos\quad\Longleftrightarrow\quad F\in\sos,
\]
because the multiplier is the invertible square
\((\sqrt s\otimes\sqrt t)^2\).

If \(X=BB^T\) over \(K_c\), \(Y=CC^T\) over \(K_t\), then
\begin{equation}
 \sum_{i,j}X_{ij}\otimes Y_{ij}
 =\sum_{a,b}\left(\sum_iB_{ia}\otimes C_{ib}\right)^2.
 \label{eq:pair}
\end{equation}
Singular PSD matrices are allowed.

For nonnegative source weights \(a_j\) and testing weights \(w_i\)
with \(\sum_iw_i=1\), write \(f=\sum_ja_jf_j\), \(q=\sum_iw_iq_i\).
If \([f_j,q_i]=\sum_k h_{ijk}^2\), then
\begin{equation}
 [f,q]=\sum_{i,j,k}
 \left((\sqrt{a_j}\otimes\sqrt{w_i})h_{ijk}\right)^2.
 \label{eq:mixture}
\end{equation}
This is finite linearity in \(f\) and affineness in \(q\).
The testing-weight sum preserves the affine constant. Zero weights
are never inverted. \src{O, FA10--FA12}

\subsection{Positive residue gives an exact tensor SOS}

Let \(\mathcal O_c,\mathcal O_t\) be the rings of elements bounded by
real constants and set \(\mathcal B=\mathcal O_c\otimes_\R\mathcal O_t\).
Their standard-part maps induce a real residue map on \(\mathcal B\).

\begin{lemma}[Positive residue]
Every \(a\in\mathcal B\) with \(\res(a)>0\) is a finite sum of
positive pure tensors, hence an SOS with rank-one roots.
\end{lemma}

\begin{proof}
Expand bounded factors around their real residues:
\[
 a=\epsilon+\sum_{j=1}^{N}x_j\otimes y_j,\qquad\epsilon>0,
\]
where each term has an infinitesimal factor. Split
\(\epsilon=\sum_j\epsilon_j\) into positive real reserves.
If \(x_j\) is infinitesimal, choose real \(M_j>|y_j|\). Exactly,
\[
 \epsilon_j+x_j\otimes y_j
 =\frac12\left[
 (\epsilon_j/M_j+x_j)\otimes(M_j+y_j)
 +(\epsilon_j/M_j-x_j)\otimes(M_j-y_j)\right].
\]
All four factors are positive in their own fields, so take their
separate square roots. Exchange factors when the infinitesimal is
\(y_j\). The case \(N=0\) is immediate.
\end{proof}

\begin{lemma}[Strict bounded Gram]
If \(H\in\operatorname{Sym}_r(\mathcal B)\) has positive-definite real
residue, then \(v^THv\in\sos\) for every \(v\in\T^r\).
The rows \(v_i\) need not be bounded.
\end{lemma}

\begin{proof}
A real congruence makes the residue \(I\). Choose real \(\delta>0\)
with \((r-1)\delta<1\). The exact identity is
\[
\begin{aligned}
 H={}&\sum_i[H_{ii}-(r-1)\delta]e_ie_i^T\\
 &+\sum_{i<j}\frac{\delta+H_{ij}}2(e_i+e_j)(e_i+e_j)^T\\
 &+\sum_{i<j}\frac{\delta-H_{ij}}2(e_i-e_j)(e_i-e_j)^T.
\end{aligned}
\]
Every scalar coefficient has positive real residue. Apply the preceding
lemma and evaluate on \(v\).
\end{proof}

There is no infinite expansion or mixed square root. The hard application
step is producing an exact full-parent Gram with bounded coefficients
and \emph{strict} real residue. Positive infinitesimal eigenvalues do
not meet the latter condition. \src{G}

\subsection{Regular contact and the additional rank-two condition}

If a real polynomial \(F(s,t)\) satisfies
\(F(s,s)=0\) and \(\partial_tF(s,s)=0\), then polynomial expansion in
\(h=t-s\) gives exactly
\[
 F(s,t)=(t-s)^TH(s,t)(t-s),\qquad
 H(p,p)=\tfrac12\partial_t^2F(p,p).
\]
No higher term is discarded. When independent source and test points
specialize to \(p\) and \(H(p,p)\succ0\), strict Gram applies.

For a smooth affine variety, the analogous statement uses an exact
tensor in the square of the diagonal ideal and a positive intrinsic
quadratic class. Actual separate lifts to a common valid chart are
needed. Scheiderer's diagonal lemma supplies this mechanism under its
precise hypotheses. \src{TF; S2}

Every real linear combination of genuine divided differences
\(f_i(s)\otimes1-1\otimes f_i(t)\) is again one divided difference,
so has rank at most two. Multiplication by rank-one coefficient roots
preserves this bound. Arbitrary rows that individually have rank two
need not have a real span consisting of rank-two tensors.
The October qualitative route requires finite root rank, not the
additional uniform bound two.

\section{The exact source funnel}

For a fixed CHSH orientation,
\[
 H=1-\tfrac12(E_{00}-E_{01}+E_{10}+E_{11}),\qquad
 f=dH+4\sum_\xi \xi p_\xi,
\]
with eight nonnegative event coefficients
\(\xi=\lambda,n,m,C,A,D,\alpha,\beta\).
For \(D_u=1+u^2\), \(D_v=1+v^2\), and state \((W,X,Y,Z)\), the
complete parent is
\begin{align}
F={}&4\lambda D_u(W+vX)^2+4nD_uD_vX^2
 +4m(vW-X+uvY-uZ)^2 \notag\\
 &+4CD_v(uW-Y)^2+4AD_uD_vY^2 \notag\\
 &+(4D+2d)(uW+uvX-Y-vZ)^2 \notag\\
 &+(4\alpha+2d)D_u(vY-Z)^2
 +(4\beta+2d)D_v(X+uZ)^2-2dD_uD_vZ^2.
 \label{eq:parent}
\end{align}
Coefficient and testing expressions occupy separate factors.
A scalar restriction of \eqref{eq:parent} cannot replace a certificate
for the entire state form.

For homogeneous measurement pairs, with
\(\Delta_A=a_0^2+a_1^2>0\), \(\Delta_B=b_0^2+b_1^2>0\),
\[
 \mathcal F=(1\otimes\Delta_A\Delta_B\|\psi\|^2)[f,\Phi(z)].
\]
The factor is pure testing and positive for a nonzero state.
The homogeneous identity includes literal endpoints without dividing
by a vanished affine coordinate. \src{O, FA4--FA6; RFA}

\subsection{Reduction to events, closed faces and two contact families}

A valid affine functional either is nonnegative on the nonsignaling
polytope and is a nonnegative event sum, or is negative at a PR vertex.
That PR vertex and its eight neighboring deterministic vertices form
an affine evaluation basis, giving the canonical coefficient
representation above.

For \(d>0\), the signed \(Z=0\) restriction is
\[
 2d\{(uW+uvX-Y)^2+D_uv^2Y^2+D_vX^2\}\succ0
 \quad(uv\ne0).
\]
Thus the three-state principal block is positive definite.
The full state determinant is
\[
 \det P=16D_u^2D_v^2u^2v^2L(U,V),\qquad U=u^2,\ V=v^2,
\]
\[
 L(U,V)=h+\sum_{e\in\{-1,0,1\}^2\setminus\{0\}}
 q_eU^{e_1}V^{e_2},\qquad q_e\ge0.
\]
The source body prints all eight nonnegative factorizations.
At a positive stationary contact,
\(a_e=q_eU_c^{e_1}V_c^{e_2}\) obey
\[
 h=-\sum_ea_e,\qquad \sum_ea_ee=0,\qquad
 \mathcal H_L=\sum_ea_eee^T.
\]
Rank-zero and rank-one support possibilities are excluded by explicit
event-incidence arguments and negative physical tests. This gives
coercivity and strict \emph{single-field} curvature at every actual
finite zero.

Conversely, nonnegative coefficients, a finite kernel state and both
angle jets zero imply physicality: the convex logarithmic Laurent
profile has a global minimum there. Logarithms are used over \(\R\);
the resulting finite ordered-field implication transfers. None of this
inverts a tensor or promotes a source-field Hessian to strict residue.
\src{SE1--SE10}

Normalize the coefficient sum to one over \(\R\).
A coercive extreme source cannot have a strictly positive Laurent
minimum, since small opposite coefficient perturbations preserve
positivity. Thus it has a finite contact.
There are four kernel equations and two angle equations, linear in
the nine coefficients. Seven or eight positive events would permit a
nonradial sign-preserving perturbation retaining all six equations,
contradicting extremality. Hence five or six events are positive.

The compatible omission patterns reduce by actual Bell symmetries to
\[
\begin{array}{ll}
\mathrm{C1}:&
 d,m,C,A,\alpha>0,\quad\lambda=n=0,\quad
 D,\beta\ge0,\quad D+\beta>0,\\
\mathrm{C2}:&
 d,m,C,A,D,\alpha>0,\quad\lambda=\beta=0,\quad n\ge0.
\end{array}
\]
Incompatible omissions enter two closed-face types.
The compact normalized real source base has affine dimension at most
eight, so nine extreme sources suffice. This fixed finite assertion
then transfers to every source field. \src{SE11--SE15}

\subsection{The closed faces have explicit lifts}

For a symmetric quadratic matrix polynomial \(K(t)=K_0+tK_1+t^2K_2\),
the retained argument proves
\[
 K(t)\succeq0\ \forall t
 \ \Longleftrightarrow\
 \exists S=-S^T:\quad
 \begin{pmatrix}K_0&K_1/2+S\\K_1/2-S&K_2\end{pmatrix}\succeq0.
\]
The proof uses a closed Gram cone, a finite matrix-moment decomposition
and separation, including singular cases.
It yields one \(6\times6\) block for \(\lambda=m=0\), and two
\(3\times3\) blocks for \(\lambda=\alpha=0\).

Both have displayed strictly feasible points. Strict feasibility bounds
the trace of PSD witnesses when their adjoint images stay bounded, so
the adjoint image is closed and the evaluation witness is attained.
The finite witness statement transfers to the testing field, and
\eqref{eq:pair} gives tensor squares. This is an actual analytical
closed-face input. \src{CF1--CF11; RFA}

\section{Finite closure and existence}

The remaining local proposition must certify the complete parent for
every actual physical C1/C2 source, every independent projective test,
every state, all allowed literal zero coefficients and all positive
openings. The initial scale
\[
 \delta_{\rm tot}=
 \frac{d}{d+\lambda+n+m+C+A+D+\alpha+\beta}
\]
is split before a \(d=1\) gauge. Positive-unit and positive-infinitesimal
total scale are distinct cases; \(d=1\) does not itself imply bounded
event coefficients.

Assuming this full local proposition, each event source has a direct
projector certificate, each closed-face source uses matrix pairing,
and each C1/C2 source uses its local parent identity.
If \((s\otimes1)\mathcal F=\sum_kh_k^2\), \(s>0\), restoration gives
\[
 [f,\Phi(z)]=
 \sum_k\left[
 \left(s^{-1/2}\otimes
       (\Delta_A\Delta_B\|\psi\|^2)^{-1/2}\right)h_k
 \right]^2.
\]
All inverses are pure field operations.
The finite source decomposition and nine-atom testing decomposition,
combined with \eqref{eq:mixture}, give \eqref{eq:universal}.
Taking two copies of one arbitrary real closed field then yields
\eqref{eq:lift}. \src{O; LC9--LC10}

This final implication is sound with those hypotheses.
E1--E3 ask for an evaluated common feature list, a concrete affine
pencil and projection, and both-inclusion validation of its data.
They are downstream constructive outputs, not extra hypotheses of
the qualitative criterion.

